% TOP_PROBLEM: {"id":30007513,"problem_number":"LOCAL-30007513","title":"Sectoral shock propagation","category_id":21,"category":{"id":21,"name":"economics","display_name":"Economics","description":"Open economic theory statements, published research agendas, and proposed scoped specifications; question provenance and historical priority are qualified per record.","slug":"economics","order_index":21,"created_at":"2026-10-08T00:00:00Z"},"status":"provenance_pending","external_url":null,"legacy_ids":[],"published":false,"statement":"Characterize how a measured production network and collateral constraints change the size of a sectoral shock required to trigger a specified aggregate recession.","economics":{"original_id":2,"field":"Business cycles and macro dynamics","kind":"theoretical","formalization_basis":"proposed_specification","source_status":"not_verified","priority_status":"not_established","priority_note":"No explicit primary open-problem statement matching this catalogue item was verified; supporting research is not evidence of first formulation.","earliest_verified_reference":null,"current_reference":{"authors":["Jorge Miranda-Pinto","Eugenio Rojas","Felipe Saffie","Alvaro Silva"],"title":"Connected for Better or Worse? The Role of Production Networks in Financial Crises","year":2026,"url":"https://www.bostonfed.org/-/media/Documents/Workingpapers/PDF/2026/WP2601.pdf","doi":"10.29412/res.wp.2026.01","locator":"Abstract and Section 1, pp. 1–3 (December 2025 manuscript issued as Working Paper 26-1)","evidence_summary":"The paper derives a network amplification statistic with collateral constraints and provides cross-country evidence for Sudden Stops. This is a substantive model-specific answer, not an explicit statement that a general network-propagation problem remains open."},"formalization_note":"New model-class specification. The cited paper already solves important special cases; no exact general open statement was verified."},"statusQualification":"Provenance pending: an explicit published statement of this exact open question was not verified. This is a catalogue-authored candidate specification, not a certified open conjecture. New model-class specification. The cited paper already solves important special cases; no exact general open statement was verified.","statusReviewedAt":"2026-10-08"}
% FIRST_POSTED: 2026-10-08
% ENTRY_KIND: statement_only
% !TeX program = lualatex
\documentclass[11pt]{article}
\usepackage{fontspec}
\usepackage[a4paper,margin=25mm]{geometry}
\usepackage{amsmath,amssymb,mathtools}
\usepackage{xurl}
\usepackage[hidelinks]{hyperref}
\newcommand{\sref}[2]{\hyperlink{source:#1:#2}{[S#2]}}
\setlength{\emergencystretch}{3em}
\begin{document}
% problemId:30007513
\section*{Sectoral shock propagation}\label{30007513}
\subsection{Definitions and mathematical statement}
Consider a finite economy with sectors \(i=1,\ldots,n\). Sector \(i\) combines labor and intermediate inputs using a constant-returns production function, has productivity \(A_i\), and purchases an expenditure share \(\omega_{ij}\) from sector \(j\). Households have specified preferences and hold debt subject to a collateral constraint whose collateral is a declared function of sectoral profits. Fix the input-share matrix \(\Omega=(\omega_{ij})\), initial balance sheets, substitution elasticities, and a rule for monetary and fiscal policy. Assume an equilibrium exists and specify an equilibrium selection if constraints generate multiplicity.
For a productivity loss \(s>0\) in sector \(j\), let \(Y_j(s)\) be equilibrium real aggregate output and define
\[
R_j(s)=1-\frac{Y_j(s)}{Y_j(0)},\qquad
s_j(q)=\inf\{s>0:R_j(s)\ge q\},
\]
where \(Y_j(0)>0\), and \(q\in(0,1)\) is a predetermined recession threshold and the infimum is \(+\infty\) if the set is empty. Characterize and quantitatively identify how network structure changes \(s_j(q)\) when financial conditions and shock sizes are held fixed. Separate first-order production effects from nonlinear collateral amplification; establish which conclusions survive plausible substitutions and alternative collateral definitions. Validate amplification predictions against independently identified sectoral shocks. Existing sufficient-statistic results for particular models are admissible benchmarks, not evidence that this entire class is unsolved.

\par\textbf{Formulation and scope.} New model-class specification. The cited paper already solves important special cases; no exact general open statement was verified.

\par\textbf{Open-question provenance.} An explicit published open-question statement was not verified. Sources below are supporting literature and must not be treated as a first listing.

\par\textbf{Historical priority.} No explicit primary open-problem statement matching this catalogue item was verified; supporting research is not evidence of first formulation.

\subsection{Short English statement}
Characterize how a measured production network and collateral constraints change the size of a sectoral shock required to trigger a specified aggregate recession.

\subsection{Sources}
\begin{itemize}\raggedright
\item[\textnormal{[S1]}]\hypertarget{source:30007513:1}{} Jorge Miranda-Pinto, Eugenio Rojas, Felipe Saffie, Alvaro Silva. 2026. Connected for Better or Worse? The Role of Production Networks in Financial Crises. Abstract and Section 1, pp. 1–3 (December 2025 manuscript issued as Working Paper 26-1).
\url{https://www.bostonfed.org/-/media/Documents/Workingpapers/PDF/2026/WP2601.pdf}.
DOI: 10.29412/res.wp.2026.01.
{\small\textit{Supporting or current literature; see evidence qualification. The paper derives a network amplification statistic with collateral constraints and provides cross-country evidence for Sudden Stops. This is a substantive model-specific answer, not an explicit statement that a general network-propagation problem remains open.}}
\end{itemize}
\end{document}
